% MÓVIL - diseño de 27 casos (CP-055..CP-081)
% GUÍA EQUIPO: Responsable asignado: Carlos Ramos (Móvil).
\paragraph{Módulo Móvil --- Seguridad y operación de entregas (27 casos: CP-055 a CP-081).}

\begin{disenio}
\caso{CP-055}{\textbf{Acceso sin sesión JWT a rutas, pedidos y configuración.} NR-01, NR-03, NR-07, NR-23; NP-14 (C06.1, Alta)}{Solicitudes HTTP sin cabecera Authorization a GET /\allowbreak{}api/\allowbreak{}v1/\allowbreak{}drivers/\allowbreak{}me/\allowbreak{}routes, GET /\allowbreak{}api/\allowbreak{}v1/\allowbreak{}shop/\allowbreak{}orders y GET /\allowbreak{}api/\allowbreak{}v1/\allowbreak{}admin/\allowbreak{}shop-settings.}{1) Invocar los tres endpoints sin token JWT. 2) Inspeccionar códigos HTTP y cuerpos de respuesta.}{Todas las peticiones se rechazan con HTTP 401 sin devolver rutas, pedidos ni parámetros de configuración.}{100 \% de los endpoints protegidos responden 401 y sus respuestas no filtran datos internos ni listados.}
\caso{CP-056}{\textbf{Aislamiento entre CLIENT-B y pedido/\allowbreak{}comprobante de CLIENT-A.} NR-01, NR-03, NR-07, NR-23; NP-14 (C06.2, Alta)}{JWT de CLIENT-B y UUID de pedido/\allowbreak{}comprobante perteneciente a CLIENT-A.}{1) Consultar GET /\allowbreak{}api/\allowbreak{}v1/\allowbreak{}shop/\allowbreak{}orders/\allowbreak{}:\allowbreak{}id y /\allowbreak{}receipt como CLIENT-B. 2) Intentar mutar o consultar /\allowbreak{}orders/\allowbreak{}:\allowbreak{}id ajeno.}{El backend rechaza cada intento con HTTP 403 o 404 sin exponer montos, destinatario ni comprobante de CLIENT-A.}{Ninguna respuesta revela datos de CLIENT-A y el pedido original permanece sin cambios.}
\caso{CP-057}{\textbf{Aislamiento de DRIVER-B ante ruta y parada de DRIVER-A.} NR-01, NR-03, NR-07, NR-23; NP-14 (C06.3, Alta)}{Sesión JWT de DRIVER-B y routeId/\allowbreak{}stopId pertenecientes a DRIVER-A.}{1) Abrir Mi Ruta en móvil con DRIVER-B. 2) Consultar ruta ajena e intentar iniciar, fallar o cerrar parada/\allowbreak{}ruta de DRIVER-A.}{Móvil no lista rutas ajenas y la API rechaza consultas y mutaciones sobre DRIVER-A con HTTP 403 o 404.}{Ruta y paradas de DRIVER-A conservan su estado intacto sin fugas de información a DRIVER-B.}
\caso{CP-058}{\textbf{Registro público con rol ADMIN rechazado sin escalación.} NR-01, NR-03, NR-07, NR-23; NP-14 (C06.4, Alta)}{Payload de registro con name, email nuevo válido, password válido y campo adicional role: ADMIN.}{1) Enviar POST /\allowbreak{}api/\allowbreak{}v1/\allowbreak{}auth/\allowbreak{}register con role ADMIN. 2) Intentar iniciar sesión con esas credenciales y verificar usuarios.}{El servidor rechaza la petición con HTTP 400 o 403 y no crea cuenta ADMIN ni usuario privilegiado.}{El intento de login posterior falla y no existe ninguna cuenta registrada con el correo enviado.}
\caso{CP-059}{\textbf{Intento de CLIENT de modificar configuración o cobertura.} NR-01, NR-03, NR-07, NR-23; NP-14 (C06.5, Alta)}{JWT de CLIENT-A con payload válido para PATCH /\allowbreak{}api/\allowbreak{}v1/\allowbreak{}admin/\allowbreak{}shop-settings y cambio en /\allowbreak{}api/\allowbreak{}v1/\allowbreak{}admin/\allowbreak{}store-coverage.}{1) Leer configuración actual como ADMIN. 2) Enviar mutaciones de settings y cobertura como CLIENT-A. 3) Releer configuración.}{Ambas peticiones de CLIENT-A se rechazan con HTTP 403 y los valores de tienda y cobertura no cambian.}{Comparación antes/\allowbreak{}después como ADMIN confirma configuración y polígono exactamente iguales.}
\caso{CP-060}{\textbf{Acceso autorizado por titular, ADMIN y DRIVER sobre recursos propios.} NR-01, NR-03, NR-07, NR-23; NP-14 (C06.6, Alta)}{JWT válidos de CLIENT-A (pedido propio), ADMIN (configuración/\allowbreak{}despacho) y DRIVER-A (ruta propia).}{1) CLIENT-A lee su pedido/\allowbreak{}comprobante. 2) ADMIN consulta settings y despacho. 3) DRIVER-A consulta sus rutas en móvil/\allowbreak{}API.}{Cada actor accede con HTTP 200 únicamente a sus recursos permitidos sin exposición de datos de terceros.}{Respuestas corresponden exactamente a los recursos autorizados de cada rol según contrato.}
\caso{CP-061}{\textbf{Pedido de hoy elegible en propuesta, ruta y asignación.} NR-16, NR-20; NP-07, NP-13 (C11.1, Alta)}{Pedido con scheduledDeliveryDate del día operativo actual en America/\allowbreak{}Guayaquil y DRIVER-A disponible.}{1) Consultar GET /\allowbreak{}api/\allowbreak{}v1/\allowbreak{}route-proposals. 2) Crear ruta con POST /\allowbreak{}api/\allowbreak{}v1/\allowbreak{}routes. 3) Asignar DRIVER-A con POST /\allowbreak{}routes/\allowbreak{}:\allowbreak{}id/\allowbreak{}assign.}{La propuesta incluye el pedido con isToday=\allowbreak{}true, permite crear borrador y asignar la ruta a DRIVER-A.}{Propuesta, ruta DRAFT/\allowbreak{}ASSIGNED y asignación quedan persistidas para la fecha operativa actual.}
\caso{CP-062}{\textbf{Pedido futuro admitido en compra y bloqueado en ruta/\allowbreak{}assign.} NR-16, NR-20; NP-07, NP-13 (C11.2, Alta)}{Pedido válido con fecha programada de mañana en POST /\allowbreak{}api/\allowbreak{}v1/\allowbreak{}shop/\allowbreak{}checkout y POST /\allowbreak{}api/\allowbreak{}v1/\allowbreak{}orders/\allowbreak{}admin.}{1) Crear pedidos futuros. 2) Consultar GET /\allowbreak{}api/\allowbreak{}v1/\allowbreak{}route-proposals?date=\allowbreak{}mañana. 3) Intentar POST /\allowbreak{}routes y /\allowbreak{}assign para mañana.}{Las compras futuras se crean (201) y permiten preview, pero POST /\allowbreak{}routes y assign se rechazan con 409 ROUTE\_DATE\_NOT\_TODAY.}{No se crean rutas ni asignaciones futuras; pedidos permanecen pendientes para su fecha programada.}
\caso{CP-063}{\textbf{Ausencia registrada hoy impide asignar repartidor.} NR-16, NR-20; NP-07, NP-13 (C11.3, Alta)}{Registro en POST /\allowbreak{}api/\allowbreak{}v1/\allowbreak{}drivers/\allowbreak{}:\allowbreak{}driverId/\allowbreak{}absences con startDate/\allowbreak{}endDate de hoy y ruta DRAFT lista.}{1) Registrar ausencia de hoy para DRIVER-B. 2) Consultar disponibilidad. 3) Intentar asignar DRIVER-B en POST /\allowbreak{}routes/\allowbreak{}:\allowbreak{}id/\allowbreak{}assign.}{Disponibilidad reporta isAvailable=\allowbreak{}false y el intento de asignación se rechaza sin asignar la ruta.}{La ruta permanece sin conductor (DRAFT) y no se genera reserva ni salida para DRIVER-B.}
\caso{CP-064}{\textbf{Ausencia futura no altera ni reasigna rutas activas de hoy.} NR-16, NR-20; NP-07, NP-13 (C11.4, Alta)}{Ausencia en POST /\allowbreak{}api/\allowbreak{}v1/\allowbreak{}drivers/\allowbreak{}:\allowbreak{}driverId/\allowbreak{}absences con fechas futuras (mañana o posterior) para DRIVER-A.}{1) Consultar ruta activa de hoy de DRIVER-A. 2) Registrar ausencia futura para DRIVER-A. 3) Reconsultar ruta y paradas de hoy.}{La ausencia futura queda registrada, y la ruta asignada o iniciada de hoy mantiene mismo conductor y paradas.}{No ocurre reasignación automática, cancelación ni alteración de reservas en la ruta de hoy.}
\caso{CP-065}{\textbf{Pedido urgente con ruta IN\_PROGRESS no altera ruta iniciada.} NR-20; NP-09 (C15.1, Alta)}{Pedido nuevo del mismo día registrado tras iniciar R1; consulta a GET /\allowbreak{}api/\allowbreak{}v1/\allowbreak{}dispatch/\allowbreak{}orders/\allowbreak{}:\allowbreak{}id/\allowbreak{}proposal.}{1) Registrar snapshot de R1. 2) Crear pedido urgente para hoy. 3) Consultar propuesta de urgencia e inbox; reconsultar R1.}{La propuesta sugiere NEW\_ROUTE o BLOCKED sin proponer inserción en R1; R1 permanece intacta.}{Paradas, orden y reserva de R1 son idénticos antes y después; no hay inserción automática.}
\caso{CP-066}{\textbf{Decisión manual ADMIN en urgencias: INSERT con cupo o NEW\_ROUTE.} NR-20; NP-09 (C15.2, Alta)}{Pedido urgente de hoy, respuesta de GET /\allowbreak{}dispatch/\allowbreak{}orders/\allowbreak{}:\allowbreak{}id/\allowbreak{}proposal y confirmación explícita en POST /\allowbreak{}dispatch/\allowbreak{}interventions.}{1) Consultar propuesta de urgencia. 2) Confirmar manualmente INSERT (cupo \ensuremath{<}4) o NEW\_ROUTE. 3) Verificar rutas e historial auditado.}{Solo tras confirmación explícita ADMIN se inserta en la ruta elegible (\ensuremath{<}4 paradas) o se crea nueva ruta con auditoría.}{Ninguna ruta supera 4 paradas, el cambio queda auditado y rutas no seleccionadas permanecen intactas.}
\caso{CP-067}{\textbf{Entrega fallida en móvil conserva parada e historial pendiente.} NR-23, NR-24; NP-10 (C18.1, Alta)}{Token DRIVER-A, stopId pendiente y motivo RECIPIENT\_ABSENT o RECIPIENT\_REFUSED en POST /\allowbreak{}api/\allowbreak{}v1/\allowbreak{}routes/\allowbreak{}:\allowbreak{}id/\allowbreak{}stops/\allowbreak{}:\allowbreak{}stopId/\allowbreak{}failure.}{1) Registrar entrega fallida desde móvil con motivo válido. 2) Consultar detalle de ruta, parada y seguimiento.}{Parada pasa a FAILED conservando intento/\allowbreak{}motivo en historial y el pedido queda en estado pendiente de gestión.}{No marca DELIVERED ni CANCELLED, conserva la parada original y el intento auditado.}
\caso{CP-068}{\textbf{Registro de fallo por actor o conductor no autorizado rechazado.} NR-23, NR-24; NP-10 (C18.2, Alta)}{Tokens de DRIVER-B y CLIENT-A sobre POST /\allowbreak{}api/\allowbreak{}v1/\allowbreak{}routes/\allowbreak{}:\allowbreak{}id/\allowbreak{}stops/\allowbreak{}:\allowbreak{}stopId/\allowbreak{}failure de DRIVER-A con motivo válido.}{1) Enviar solicitud de fallo como DRIVER-B y como CLIENT-A. 2) Consultar parada y pedido después de cada intento.}{Cada intento no autorizado se rechaza con HTTP 403 o 409 sin registrar intento fallido.}{La parada sigue PENDING y el pedido no muestra transición ni historial modificado.}
\caso{CP-069}{\textbf{Gestión post-fallo ADMIN: CONTACT y RESCHEDULE con nueva fecha.} NR-23, NR-24; NP-10 (C18.3, Alta)}{Acciones CONTACT y RESCHEDULE en POST /\allowbreak{}api/\allowbreak{}v1/\allowbreak{}dispatch/\allowbreak{}interventions con nueva fecha/\allowbreak{}hora acordada válida.}{1) Registrar CONTACT tras contacto al cliente. 2) Ejecutar RESCHEDULE con fecha/\allowbreak{}hora nueva. 3) Reintentar o consultar historial.}{El pedido queda reprogramado con la nueva fecha/\allowbreak{}hora y conserva el historial de fallo previo en la ruta original.}{Sin duplicar asignación activa ni borrar el intento previo; nueva fecha queda lista para planificación.}
\caso{CP-070}{\textbf{Gestión post-fallo ADMIN: acuerdo de retiro en tienda (PICKUP\_AGREE).} NR-23, NR-24; NP-10 (C18.4, Alta)}{Payload con action PICKUP\_AGREE en POST /\allowbreak{}api/\allowbreak{}v1/\allowbreak{}dispatch/\allowbreak{}interventions para el pedido fallido.}{1) Enviar PICKUP\_AGREE como ADMIN. 2) Consultar seguimiento del pedido como cliente y en despacho.}{El pedido pasa a estado pendiente de retiro en tienda sin marcarlo aún como entregado ni reponer stock.}{Seguimiento muestra pendiente de retiro y conserva registro histórico de la visita fallida.}
\caso{CP-071}{\textbf{Confirmación de retiro (PICKUP\_CONFIRM) y rechazo de repetición.} NR-23, NR-24; NP-10 (C18.5, Alta)}{Acción PICKUP\_CONFIRM en POST /\allowbreak{}api/\allowbreak{}v1/\allowbreak{}dispatch/\allowbreak{}interventions y repetición inmediata de PICKUP\_CONFIRM.}{1) Confirmar retiro físico con PICKUP\_CONFIRM. 2) Repetir PICKUP\_CONFIRM sobre el mismo pedido. 3) Consultar estado e historial.}{La primera confirmación marca entrega por retiro; la repetición se rechaza con 409 PICKUP\_ALREADY\_CONFIRMED.}{Existe una sola entrega por retiro sin registros duplicados ni movimientos adicionales de inventario.}
\caso{CP-072}{\textbf{Recorrido E2E: creación de compra y verificación de estado inicial.} NR-01, NR-05, NR-07, NR-19--NR-21, NR-23; Transversal (C25.1, Alta)}{Carrito válido de CLIENT-A con fecha de hoy, dirección cubierta y medio de pago demo aprobado.}{1) Completar checkout en Shop. 2) Abrir historial/\allowbreak{}seguimiento del cliente y consultar pedido y comprobante.}{Se crea pedido en estado inicial pendiente y comprobante congelado visibles en el seguimiento de CLIENT-A.}{Importes, stock descontado, comprobante y estado inicial coinciden exactamente con el backend.}
\caso{CP-073}{\textbf{Recorrido E2E: propuesta, confirmación y asignación en día operativo.} NR-01, NR-05, NR-07, NR-19--NR-21, NR-23; Transversal (C25.2, Alta)}{Sesión ADMIN, pedido de hoy elegible y DRIVER-A disponible sin ausencias ni rutas activas.}{1) Generar propuesta en web/\allowbreak{}API. 2) Confirmar borrador (DRAFT). 3) Asignar ruta a DRIVER-A y verificar en app móvil.}{Borrador se persiste, queda asignado (ASSIGNED) a DRIVER-A y aparece en Mi Ruta de la app móvil.}{Pedido, parada, reserva del conductor y lista en móvil están sincronizados sin duplicados.}
\caso{CP-074}{\textbf{Recorrido E2E: progreso móvil reflejado en seguimiento del cliente.} NR-01, NR-05, NR-07, NR-19--NR-21, NR-23; Transversal (C25.3, Alta)}{Sesión DRIVER-A en app móvil y sesión CLIENT-A en seguimiento web del mismo pedido.}{1) DRIVER-A inicia ruta en móvil. 2) Marca parada en camino y luego entregada. 3) CLIENT-A refresca seguimiento.}{Cada transición persistida (IN\_PROGRESS, IN\_TRANSIT, DELIVERED) se refleja en tiempo real o consulta en seguimiento cliente.}{Estados visibles en cliente y móvil coinciden con los persistidos en backend sin ubicaciones GPS inventadas.}
\caso{CP-075}{\textbf{Recorrido E2E: desconexión móvil y recuperación al reconectar.} NR-01, NR-05, NR-07, NR-19--NR-21, NR-23; Transversal (C25.4, Alta)}{Desconexión de red controlada en móvil antes/\allowbreak{}durante consulta o acción de parada y reconexión posterior.}{1) Cortar red en móvil y observar aviso ante intento de acción. 2) Reconectar red. 3) Reintentar o refrescar y verificar API.}{Sin red no finge éxito; al reconectar recupera el estado real persistido sin duplicar entregas ni inventar cola offline.}{Backend registra una sola transición válida y el móvil refleja exactamente el estado del servidor.}
\caso{CP-076}{\textbf{Cierre de ruta móvil rechazado con visita aún PENDING.} NR-11, NR-20, NR-23; NP-13, NP-14 (C28.1, Alta)}{Token DRIVER-A y ruta con parada pendiente hacia POST /\allowbreak{}api/\allowbreak{}v1/\allowbreak{}routes/\allowbreak{}:\allowbreak{}id/\allowbreak{}complete con cuerpo \{\}.}{1) Verificar que existe parada PENDING. 2) Intentar cerrar ruta desde móvil o POST /\allowbreak{}routes/\allowbreak{}:\allowbreak{}id/\allowbreak{}complete.}{El servidor rechaza el cierre con HTTP 409 y la ruta continúa en estado IN\_PROGRESS.}{No se registra completedAt, la ruta no pasa a COMPLETED y el conductor sigue ocupado.}
\caso{CP-077}{\textbf{Cierre de ruta COMPLETED con visitas resueltas y una FAILED.} NR-11, NR-20, NR-23; NP-13, NP-14 (C28.2, Alta)}{Token DRIVER-A y cuerpo \{\} en POST /\allowbreak{}api/\allowbreak{}v1/\allowbreak{}routes/\allowbreak{}:\allowbreak{}id/\allowbreak{}complete sobre ruta con paradas resueltas.}{1) Ejecutar cierre de ruta como DRIVER-A. 2) Consultar ruta, parada FAILED y disponibilidad del conductor.}{Ruta pasa a COMPLETED; la parada FAILED sigue fallida y el conductor aún no queda disponible para otra ruta.}{completedAt registrado, pedido FAILED sin marcar DELIVERED y disponibilidad bloqueada hasta retorno.}
\caso{CP-078}{\textbf{Bloqueo de nueva salida antes de confirmar regreso (returnedAt).} NR-11, NR-20, NR-23; NP-13, NP-14 (C28.3, Alta)}{JWT ADMIN, routeId de nueva ruta DRAFT de hoy y driverId de DRIVER-A con returnedAt IS NULL.}{1) Confirmar ruta previa COMPLETED y returnedAt nulo. 2) Intentar asignar nueva ruta a DRIVER-A en POST /\allowbreak{}routes/\allowbreak{}:\allowbreak{}id/\allowbreak{}assign.}{La asignación de nueva salida se rechaza con HTTP 409 mientras no se confirme el regreso físico a tienda.}{La segunda ruta permanece DRAFT sin conductor y no existe doble reserva activa para DRIVER-A.}
\caso{CP-079}{\textbf{Confirmación de regreso por titular libera y actualiza disponibilidad.} NR-11, NR-20, NR-23; NP-13, NP-14 (C28.4, Alta)}{Token DRIVER-A y cuerpo \{\} en POST /\allowbreak{}api/\allowbreak{}v1/\allowbreak{}routes/\allowbreak{}:\allowbreak{}id/\allowbreak{}return (probando también con ausencia activa para hoy).}{1) Confirmar regreso como DRIVER-A en móvil/\allowbreak{}API. 2) Consultar ruta, returnedAt y disponibilidad efectiva (con/\allowbreak{}sin ausencia).}{Registra returnedAt, libera reserva de ruta y actualiza disponibilidad; si existe ausencia hoy, sigue no disponible.}{returnedAt persistido, reserva liberada y disponibilidad efectiva respeta ausencias vigentes sin habilitar indebidamente.}
\caso{CP-080}{\textbf{Repetición de regreso a tienda sin duplicar ni liberar dos veces.} NR-11, NR-20, NR-23; NP-13, NP-14 (C28.5, Alta)}{Token DRIVER-A y segunda invocación con cuerpo \{\} a POST /\allowbreak{}api/\allowbreak{}v1/\allowbreak{}routes/\allowbreak{}:\allowbreak{}id/\allowbreak{}return sobre la misma ruta.}{1) Repetir POST /\allowbreak{}api/\allowbreak{}v1/\allowbreak{}routes/\allowbreak{}:\allowbreak{}id/\allowbreak{}return con DRIVER-A. 2) Comparar returnedAt, auditoría y estado del conductor.}{No duplica registro de regreso ni ejecuta doble liberación; responde idempotente o conflicto controlado según contrato.}{returnedAt original se mantiene sin cambios y la disponibilidad/\allowbreak{}reservas del conductor no sufren alteración.}
\caso{CP-081}{\textbf{Intento de cierre o regreso por DRIVER-B rechazado sobre ruta ajena.} NR-11, NR-20, NR-23; NP-13, NP-14 (C28.6, Alta)}{Token de DRIVER-B hacia POST /\allowbreak{}api/\allowbreak{}v1/\allowbreak{}routes/\allowbreak{}:\allowbreak{}id/\allowbreak{}complete y POST /\allowbreak{}api/\allowbreak{}v1/\allowbreak{}routes/\allowbreak{}:\allowbreak{}id/\allowbreak{}return de DRIVER-A.}{1) Invocar /\allowbreak{}complete y /\allowbreak{}return de ruta de DRIVER-A con JWT de DRIVER-B. 2) Consultar ruta y disponibilidad de DRIVER-A.}{Ambas acciones son rechazadas con HTTP 403 o 404 sin completar ni registrar regreso en la ruta ajena.}{Estado, completedAt, returnedAt y reserva de DRIVER-A permanecen exactamente iguales.}
\end{disenio}
